% Esquisse d'un programme (cote 119, copy 1): archive cote 119, site page 2 (PDF page 3).
% sommaire. Draft compared: EsquisseFr.pdf (text layer):1-41.
% AI-transcribed by Claude Opus 5.5 (claude-opus-5-5, adaptive thinking, effort high), Message Batches API, batch pilot (msgbatch_0189Nhf1xB615YWnztubtLfc), 2026-11-12; not yet checked by eye.
\scanpage{119/2}
\uline{ESQUISSE D'UN PROGRAMME}

par Alexandre Grothendieck

\uline{Sommaire} :

1. Envoi.

2. Un "jeu de Légo-Teichmüller", et le groupe de Galois de $\overline{\mathbb{Q}}$ sur $\mathbb{Q}$.

3. Corps de nombres associés à un dessin d'enfant.

4. Polyèdres réguliers et corps finis.

5. Haro sur la topologie dite ``générale'', et réflexions heuristiques
vers une topologie dite ``modérée''.

6. ``Théories différentiables'' (à la Nash) et ``théories modérées''.

7. A la Poursuite des Champs.

8. Digressions de géométrie bidimensionelle.

9. Bilan d'une activité enseignante.

10. Epilogue.

Notes.

N.B. Les astéristiques (*) renvoient aux notes figurant au bas de
la même page, les renvois numérotés de (\textsuperscript{1}) à (\textsuperscript{7}\unsure{.}) aux notes (rajou-%
tées ultérieurement) réunies à la fin du rapport.

\mspage{2}
